%% =================================================================
%% Wei, Mei, Zhao, Xiao, Sun, Zhang, Guo.
%% "Nondestructively identifying the mechanical behavior of soft
%%  tissues using surface deformation with an explicit inverse
%%  approach."
%% Appl. Math. Modelling 134 (2024) 126--147.
%%
%% Reconstruction of page 15 (printed p. 140) as study notes.
%% Generated by: reproduce-multiple-pages-basic-tex (batch 13-16)
%%
%% Structure: Fig. 18 (optimization trajectory for ring) + Fig. 19
%% (explicit vs implicit ring reconstructions) + italic subsection
%% "Case 3" + one prose paragraph (cut off mid-sentence).
%% No equations, no tables.
%% =================================================================

\documentclass[11pt]{article}

\usepackage[margin=1in]{geometry}
\usepackage{amsmath, amssymb}
\usepackage{mathrsfs}
\usepackage{bm}
\usepackage{xcolor}
\usepackage{fancyhdr}

\pagestyle{fancy}
\fancyhf{}
\lhead{\small Y.~Mei et al.}
\rhead{\small Applied Mathematical Modelling 134 (2024) 126--147}
\cfoot{\small 140 $\to$ \arabic{page}}
\renewcommand{\headrulewidth}{0.4pt}

\setcounter{page}{1}

\begin{document}

% ---------- Figure 18 placeholder ----------
\noindent
\begin{center}
\fbox{\parbox{0.92\textwidth}{\centering\vspace{1em}
\textbf{[Figure 18 --- placeholder]}\\[0.8em]
Six-panel optimization trajectory for the \textbf{Case 2 bilayer
ring structure} (target: outer 0.386\,MPa, inner 0.129\,MPa) using
the explicit inverse scheme with 3 surface displacement datasets
(no noise), starting from a 4-layer initial guess. Each panel is
a circular ring with a jet colorbar.\\[0.4em]
\begin{tabular}{@{}ll@{}}
\textbf{(a) Initial guess}  & 4-layer ring, SM $\in [0.01, 0.04]$\,MPa \\
\textbf{(b) Iteration 20}   & SM $\in [0.145, 0.438]$\,MPa (overshoots outer layer) \\
\textbf{(c) Iteration 50}   & SM $\in [0.131, 0.397]$\,MPa, wavy interface emerging \\
\textbf{(d) Iteration 100}  & SM $\in [0.129, 0.393]$\,MPa \\
\textbf{(e) Iteration 1000} & SM $\in [0.131, 0.387]$\,MPa (near target) \\
\textbf{(f) Iteration 1322 (final)} & SM $\in [0.131, 0.384]$\,MPa
\end{tabular}\\[0.4em]
The outer-layer colorbar max converges to 0.386\,MPa
\emph{from above}: the trajectory overshoots early
($0.438 \to 0.397 \to 0.393 \to 0.387 \to 0.384$) and settles
slightly under the nominal 0.386 target. The ring interface
becomes increasingly wavy as iteration proceeds but the overall
bilayer structure is recovered faithfully.
\vspace{1em}}}
\end{center}
\vspace{0.3em}
\begin{center}
\textbf{Fig. 18.} The optimization process with four layers
initially utilizing the explicit inverse scheme with 3 surface
datasets (Without noise, Unit: MPa).
\end{center}

\vspace{0.8em}

% ---------- Figure 19 placeholder ----------
\noindent
\begin{center}
\fbox{\parbox{0.92\textwidth}{\centering\vspace{1em}
\textbf{[Figure 19 --- placeholder]}\\[0.8em]
Six-panel comparison of the \textbf{Explicit Inverse Method} (top
row) vs.\ \textbf{Implicit Inverse Method} (bottom row) on the
bilayer ring structure (\textbf{no noise}), with 3, 6, and 9
surface displacement datasets.\\[0.4em]
\textbf{Explicit (top row)}: all three reconstructions are
visually near-identical to the Fig.~17 target --- clean bilayer
ring with wavy inner interface. Colorbar ranges:
(a) 3 fields, SM $\in [0.131, 0.384]$\,MPa;
(b) 6 fields, SM $\in [0.131, 0.384]$\,MPa;
(c) 9 fields, SM $\in [0.131, 0.384]$\,MPa.
The explicit method reaches the same quality with 3 fields as
with 9 --- consistent with its low parameter count
(56 variables per the paper's earlier claim).\\[0.4em]
\textbf{Implicit (bottom row)}: reconstructions are visibly
\textbf{dominated by green/yellow artifacts} with no clear bilayer
structure recovered. Colorbar ranges show a \emph{monotonically
increasing max} as more datasets are added:
(d) 3 fields, SM $\in [0.010, 0.478]$\,MPa;
(e) 6 fields, SM $\in [0.010, 0.632]$\,MPa;
(f) 9 fields, SM $\in [0.010, 0.813]$\,MPa.
More data makes the implicit result \emph{worse}, not better ---
the 9-field max overshoots the 0.386 nominal outer-layer value by
\textbf{110\,\%}.\\[0.3em]
\textbf{Takeaway}: this is the most dramatic failure mode of the
implicit method in the entire paper. On the ring structure, with
4030 optimization variables to the explicit method's 56, the
implicit method cannot find the bilayer structure at all ---
increasing displacement data only makes it diverge more. This
suggests the implicit method's high dimensionality is
\emph{actively harmful} on this problem, not just inefficient.
\vspace{1em}}}
\end{center}
\vspace{0.3em}
\begin{center}
\textbf{Fig. 19.} The reconstructed results with varying 3, 6 and
9 surface displacement datasets (Without noise, Unit: MPa).
\end{center}

\vspace{1em}

% ---------- Case 3 subsection heading (italic, matching paper style) ----------
\noindent\textit{Case 3: A composite with inclusions embedded}

\vspace{0.5em}

\noindent
In the last example, we replicated a scenario where one or two
tumors were embedded within soft tissues with a stiffness ratio
of 5:1, as illustrated in Fig.~21. To simulate the forward problem,
we applied an ``indentation'' force on the top surface of a square
domain, discretized using 7200 triangular elements. The center
point was fixed and other points where on the bottom edge are
constrained in the vertical direction, and we only utilized the
displacements on the top surface to tackle the inverse problem.
To initiate the inverse problem, we employed four inclusions as
the initial guess (refer to Fig.~22(a)). These inclusions were
predefined in the problem

% [page ends here — continues on page 16]

\end{document}
